\begin{theorem}[The obstruction for coprime multiplicities]\hypertarget{tag-02RJ}{}\label{night:P2.9}\label{stmt:P2.9}
\StatementRecord{night:P2.9}{proved}{02RJ}{the definitions and constructions stated in the proof}

Every admissible logarithmic datum in the coprime $(3k,4l)$ family of Viaclovsky's logarithmic construction has no proper Hausdorff analytic filling at parameter zero after any finite cover.
\end{theorem}
\begin{proof}
\begin{enumerate}[label=\arabic*.,leftmargin=*]
\item Put $m_1=3k$, $m_2=4l$, $N=m_1m_2=12kl$, with $\gcd(m_1,m_2)=1$. Let $l_j$ be the $\widehat\gamma$ coefficient of the invariant finite twist. Admissibility is $\gcd(l_j,m_j)=1$. Retain $l_0$ and define $b=m_2l_1+m_1l_2$, $a=b-Nl_0$. On $w=r^{N/2}$, $c=-Nv$, and the gamma period is $(6\mu,C-s-Nv)$. The descent is
\[
 H=-\frac{m_2l_1\log t+m_1l_2\log(t-1)}{2\pi i}=bs+U,
 \qquad U=-\frac{m_1l_2\log(1-q)}{2\pi i}.
\]
The finite actions are unchanged by $c\mapsto c-N$, since $N/m_j$ is integral. The collar coefficient is $a$; adding $a\gamma$ gives a unit action on cusp charts. The rank-one filling has $N$ components with rotations $m_2l_1,m_1l_2,-b$. Admissibility makes the finite actions free on the cycle.

\item On $q=Q^N$ the origins are
\begin{align*}
 o&=\binom{-6b\mu s/N}{bs^2/(2N)-bCs/N+l_0s/2+v(bs+U)},\\
 \sigma&=-l_0z_\gamma+\binom{6a\mu v}{vU+aCv-Nav^2/2+Nav/2},
\end{align*}
with $z_\gamma=(6s\mu,-s(s+1)/2+s(C-Nv))$. Their difference is
\[
 aF_N=a\binom{6\mu(s+Nv)/N}{C(s+Nv)/N-(s+Nv)^2/(2N)+Nv/2}.
\]
On $q=Q^{Nm}$, $r=R^n$,
\begin{align*}
 \Delta_Q(aF_N)&=am\gamma-\tfrac N2am^2\delta,\\
 \Delta_R(aF_N)&=an\gamma-\tfrac N2an(n-1)\delta.
\end{align*}
All coefficients are integral because $12\mid N$. Expanding the two quadratic expressions gives the origin subtraction, the four descent identities and both covered differences.

\item Pulling the order-three and order-four tangent root charts through $k$ and $l$ powers changes their vertical determinant orders from $2,1$ to $2k,l$. The vertical determinant is therefore $f^*\OO(-1)\otimes\OO(2kS_1+lS_2)$. The determinant of the image of $df$ is $f^*\OO(2)\otimes\OO(-(3k-1)S_1-(4l-1)S_2)$; the cusp singular ideal has codimension two. Hence
\begin{align*}
 K_X^{-1}&=f^*\OO(1)\otimes\OO((1-k)S_1+(1-3l)S_2),\\
 K_X^{-N}&=f^*\OO(d),& d&=4l+3k-kl.
\end{align*}
The integer $d$ is nonzero. Otherwise $k\mid4l$, while $\gcd(k,4l)=1$, so $k=1$, giving $d=3l+3>0$.

\item Choose $K^{-N}$ when $d>0$ and $K^N$ when $d<0$; its restriction is $f^*\OO(|d|)$. On a resolved proper filling the corresponding relative direct image is coherent torsion free of rank $|d|+1$. Evaluation and graph resolution give a proper pluricanonical image. Over the field $K$ of meromorphic parameter germs, its image is a smooth rational normal curve. More explicitly, the kernels of multiplication on the direct-image modules, tensored with $K$, give its quadratic equations; on the punctured disc these are the Veronese equations and generate its ideal. The three labelled critical-value curves give distinct $K$-rational points. The coprime multiplicities are unequal, and the third point is the cusp.

\item Over an algebraic closure of $K$ there is a unique rational parameter taking those points to $0,1,\infty$. Uniqueness makes it invariant under every field automorphism over $K$, so its coefficients belong to $K$. They are meromorphic germs with finite principal parts. Composing and resolving the graph yields the marked base map required in \ref{night:P2.4}. This explicitly recovers a degree-one base parameter.

\item Modulo $\delta$, the global lattice and finite-orbit calculation are unchanged. The cross-cocycles still lie in $\Z\delta$, so \ref{night:P2.6} applies. At the corner,
\[
 A_Q=I-NmE_{32}+NmE_{41},\quad A_R=I+NnE_{41},\quad [Q,R]=-Nmnb\widehat\delta.
\]
The commuting lifts $Q_o=T_{-mb\widehat\gamma}Q$, $R_o=R$ give \eqref{eq:P2ratio}. Connectedness forces equality of the horizontal ratio zero and cusp ratio $a$. But $a\equiv m_2l_1\pmod{m_1}$ and $a\equiv m_1l_2\pmod{m_2}$, so $\gcd(a,N)=1$. This contradiction proves nonexistence. Proper pullback to the further cover preserves the punctured family and gives the implication back to every finite parameter cover.
\end{enumerate}
\end{proof}
